\section{结论}

本文将 CLIP 零样本异常定位从直接局部响应评分重新表述为图像条件化正常参照估计。局部响应可以提示潜在异常，其证据价值由当前图像的正常外观上下文决定。基于这一认识，\method{} 在给定冻结 CLIP/AnomalyCLIP 表征和 normal/abnormal prototypes 的前提下，从测试图像内部选择可靠局部证据，构造图像条件化 prototypes，并在同一语义空间内重新评分。

MVTec/VisA 受控消融支持这一机制解释：直接局部线索使用是负控，语义响应自校准是强对照，边界感知局部可靠性和保守参照校准共同形成稳定路径。本文的核心 takeaway 是一个可检验的异常证据定义：候选局部线索经过当前图像正常参照约束后，成为异常定位证据。
